% ============================================================
%  CHAPTER 10: DEPLOYMENT, LIMITATIONS, AND MAINTENANCE
% ============================================================

\chapter{DEPLOYMENT, LIMITATIONS, AND MAINTENANCE}
\label{chap:deployment}

\section{Production Deployment Architecture}
\label{sec:deployment_architecture}

Deploying QuantumSafe VPN in an enterprise or personal privacy environment involves establishing a hardened, resilient infrastructure consisting of two operational endpoints: the \textbf{Linux VPN Gateway} and the \textbf{Android Mobile Endpoints}. The gateway is deployed within an enterprise Demilitarized Zone (DMZ) or hosted on a virtual private server (VPS) with direct public WAN connectivity.

\section{Linux Gateway Server Setup and System Configuration}
\label{sec:linux_server_setup}

The gateway operates on Kali Linux 2024.x, Ubuntu Server 22.04 LTS, or Debian 12 with Linux Kernel 5.15+.

\subsection{Prerequisite System Packages}
Prior to starting the daemon, essential networking and runtime utilities are installed:
\begin{lstlisting}[language=bash, caption={Linux Server Dependency Installation.}]
sudo apt-get update && sudo apt-get install -y \
    python3 \
    python3-pip \
    python3-venv \
    iptables \
    iproute2 \
    net-tools \
    tcpdump
\end{lstlisting}

\subsection{Enabling Kernel IPv4 Packet Forwarding}
To permit the Linux kernel to route packets between the virtual \texttt{tun0} interface and the physical \texttt{eth0} WAN interface, IPv4 forwarding must be enabled:
\begin{lstlisting}[language=bash, caption={Enabling Linux Kernel IPv4 Forwarding.}]
# Enable runtime forwarding
sudo sysctl -w net.ipv4.ip_forward=1

# Persist across system reboots
echo "net.ipv4.ip_forward = 1" | sudo tee -a /etc/sysctl.d/99-vpn-forward.conf
sudo sysctl --system
\end{lstlisting}

\subsection{Configuring Firewall and iptables NAT Masquerading}
The server requires an explicit \texttt{POSTROUTING} rule in the \texttt{nat} table to translate internal VPN client IP addresses (\texttt{10.8.0.0/24}) to the server's public WAN IP:
\begin{lstlisting}[language=bash, caption={iptables NAT and Forwarding Rules.}]
WAN_IFACE="eth0"
VPN_SUBNET="10.8.0.0/24"

# Clear existing legacy rules
sudo iptables -t nat -F POSTROUTING
sudo iptables -F FORWARD

# Masquerade outbound traffic from VPN clients
sudo iptables -t nat -A POSTROUTING -s $VPN_SUBNET -o $WAN_IFACE -j MASQUERADE

# Allow bidirectional forwarding between tun0 and physical WAN
sudo iptables -A FORWARD -i tun0 -o $WAN_IFACE -j ACCEPT
sudo iptables -A FORWARD -i $WAN_IFACE -o tun0 -m state --state RELATED,ESTABLISHED -j ACCEPT

# Allow incoming TCP handshakes on VPN port 7000
sudo iptables -A INPUT -p tcp --dport 7000 -j ACCEPT
\end{lstlisting}

\section{Automated Gateway Deployment Scripting}
\label{sec:automated_deployment_scripts}

To eliminate human error during gateway initialization, the project repository provides automated management scripts in \texttt{scripts/}:
\begin{itemize}
    \item \texttt{scripts/setup\_network.sh}: Configures kernel sysctl parameters and sets up \texttt{iptables} NAT rules.
    \item \texttt{scripts/start\_server.sh}: Validates permissions, checks listening ports, and starts \texttt{vpn/tunnel\_server.py} within a detached screen session or systemd daemon.
\end{itemize}

\begin{lstlisting}[language=bash, caption={Server Launch Script (scripts/start\_server.sh).}]
#!/usr/bin/env bash
set -e

echo "[*] Initializing QuantumSafe VPN Linux Gateway..."
if [ "$EUID" -ne 0 ]; then
  echo "[-] Error: This deployment script must be run as root."
  exit 1
fi

# Execute network forwarding setup
bash "$(dirname "$0")/setup_network.sh"

# Launch Python multi-client VPN daemon
export PYTHONPATH="."
export VPN_SERVER_PORT=7000
echo "[+] Starting tunnel_server.py on port 7000..."
python3 vpn/tunnel_server.py
\end{lstlisting}

\section{Android Client Deployment and Sideloading}
\label{sec:android_client_deployment}

The Android mobile client is distributed via signed production APK artifacts:
\begin{enumerate}
    \item \textbf{Device Preparation}: On the target Android smartphone, enable ``Install unknown apps'' for the browser or file manager within \textit{Settings $\to$ Apps $\to$ Special app access}.
    \item \textbf{APK Sideloading via ADB}: Connect the device via USB with USB Debugging enabled, and execute:
    \begin{lstlisting}[language=bash, caption={ADB Sideload Installation Command.}]
adb install -r dist/QuantumSafe-VPN-release.apk
    \end{lstlisting}
    \item \textbf{Initial Launch and Configuration}:
    \begin{itemize}
        \item Open ``QuantumSafe VPN'' from the application launcher.
        \item Navigate to the \textbf{Server Config} tab and input the Linux gateway's IP address (e.g., \texttt{11.12.5.8}) and port (\texttt{7000}).
        \item Return to the \textbf{Home} tab and tap the primary \textbf{``Connect''} button.
        \item Tap \textbf{``OK''} on the Android system \texttt{VpnService} connection request dialog.
    \end{itemize}
\end{enumerate}

\section{Operational Troubleshooting and Diagnostic Runbook}
\label{sec:troubleshooting_runbook}

Table~\ref{tab:troubleshooting_guide} details common operational issues, diagnostic indicators, and verified resolution steps.

\begin{table}[htbp]
  \centering
  \small
  \caption{Operational Troubleshooting and Diagnostic Runbook.}
  \label{tab:troubleshooting_guide}
  \begin{tabular}{@{}P{3.2cm}P{4.2cm}P{5.8cm}@{}}
    \toprule
    \textbf{Observed Symptom} & \textbf{Probable Root Cause} & \textbf{Prescribed Resolution Procedure} \\
    \midrule
    \textbf{Connection Timeout (UI displays ERROR)} & TCP port 7000 blocked by intermediate firewall; server daemon inactive & 1. Verify server daemon is running: \texttt{ss -tulpn | grep 7000}.\\ & & 2. Inspect intermediate firewall / security groups.\\ & & 3. Confirm phone and server share network routing. \\
    \addlinespace
    \textbf{Tunnel Connected but No Internet Access} & Linux kernel IP forwarding disabled; \texttt{iptables} masquerading rule missing & 1. Check forwarding: \texttt{sysctl net.ipv4.ip\_forward}.\\ & & 2. Verify NAT table: \texttt{iptables -t nat -L -v -n}.\\ & & 3. Run \texttt{sudo bash scripts/setup\_network.sh}. \\
    \addlinespace
    \textbf{Signature Verification Failure} & Server identity key modified; MitM proxy active on port 7000 & 1. Inspect server log for key regeneration.\\ & & 2. Clear cached identity hash in app Settings.\\ & & 3. Re-verify server \texttt{keys/server\_mldsa65\_public.key}. \\
    \addlinespace
    \textbf{Phone UI Freezes upon Connecting} & Socket protection failure triggering recursive kernel loop & 1. Verify \texttt{protect(socket)} returned \texttt{true}.\\ & & 2. Ensure Android OS has not restricted VPN app.\\ & & 3. Restart application process. \\
    \bottomrule
  \end{tabular}
\end{table}

\section{Inherent Architectural and Technical Limitations}
\label{sec:system_limitations}

While QuantumSafe VPN delivers robust post-quantum tunneling, several technical limitations exist in the current engineering release:

\begin{enumerate}
    \item \textbf{TCP Encapsulation Overhead}: The current release tunnels IP datagrams over TCP. Under severe cellular network congestion with elevated packet loss, nested TCP retransmission timers can induce temporary throughput degradation (TCP-over-TCP meltdown).
    \item \textbf{Single-Hop Gateway Architecture}: The system operates as a client-to-gateway star topology. It does not currently support multi-hop onion routing or dynamic decentralized mesh topologies (such as Tailscale or WireGuard peer meshes).
    \item \textbf{IPv4 Exclusivity}: The virtual TUN interface captures and routes IPv4 traffic. While IPv6 is intentionally suppressed on the client to prevent DNS leaks, native IPv6 dual-stack encapsulation is not yet implemented.
    \item \textbf{Pure-Kotlin Mathematical Fallback}: While functional and memory-safe, pure-Kotlin polynomial arithmetic does not leverage platform-specific SIMD vector instructions (e.g., ARM NEON), resulting in slightly higher CPU cycle counts than highly optimized C/assembly implementations.
\end{enumerate}

\section{Security Maintenance and Dependency Lifecycle}
\label{sec:security_maintenance}

To maintain long-term security posture:
\begin{itemize}
    \item \textbf{Cryptographic Parameter Tracking}: As NIST and the cryptographic research community continue to audit finalized FIPS 203 and 204 implementations, parameter sets and polynomial round constants must be tracked against official NIST Special Publications.
    \item \textbf{Zero-Logging Enforcement}: The Linux gateway operates under strict operational hygiene, discarding all packet payload records immediately after forwarding. Session routing entries are maintained strictly in volatile RAM and purged upon socket termination.
    \item \textbf{Server Key Rotation}: Server ML-DSA-65 identity keys should be rotated annually, with public key fingerprints published out-of-band via secure administrative channels.
\end{itemize}
